\section{Discussion}

Zapit is a complete and accessible solution for head-fixed scanner-based optogenetics. The components are commercially available, and the modular system is flexible enough to accommodate a diversity of experimental needs. The \href{https://github.com/Zapit-Optostim/zapit}{Github} repository contains complete build instructions, parts lists, and visual aids: no prior expertise is required for implementation. The  software is fully documented and will be continually maintained and updated by developers and experimentalists. Known software and hardware issues are actively raised and resolved through \href{https://github.com/Zapit-Optostim/zapit/issues}{Github} issues and community discussions. 

Zapit produced robust behavioral effects across a variety of experimental protocols and brain regions with VGAT-ChR2 mice. We validated the efficacy of inactivation in this mouse strain with simultaneous electrophysiology recordings and demonstrate that the effective radius of inactivation depends on the laser power used, but was \textasciitilde{}1~mm at powers for which we observed significant behavioral perturbations. However, with finer spatial sampling and large trial numbers, it is possible to observe distinct behavioral effects at sites separated by less than 1 mm.

There are other effective approaches for rapid programmatic photostimulation in head-fixed behavioral tasks: those based on digital micro-mirror devices (DMDs \cite{Chong2020ManipulatingPerception, Russell2022All-opticalMice, Gauld2024APerception}) or spatial lights modulators (SLM). DMD-based solutions, such as the \href{https://www.mightexbio.com/polygon/}{Mightex Polygon}, provide a faster pattern-switching time but due to light loss, these systems require \textit{\textasciitilde{}100 times} more laser power than Zapit to cover the mouse dorsal cortex. SLM-based solutions are more light efficient, but devices with high resolution are very expensive and are slower than DMDs. Therefore, Zapit's scanner-based approach provides the most cost-effective and compact solution for the majority of optogenetic experiments. 

\subsection{Practical considerations when using Zapit}
\subsubsection{Surgical procedures} 
To maximize targeting flexibility and accuracy, we recommend a head-fixation device that provides an unobstructed view of the entire dorsal surface of the skull, and marking structural reference points during surgery. We further recommend a thin layer of transparent glue (\cite{Zatka-Haas2021SensoryDecision, Coen2023MouseDecisions}) or `bone' cement (200-300~$\mu m$ thick) to protect the skull as thicker layers of cement (0.5-1~mm) may reduce the efficacy of behavioral effects. For general details of \textit{in vivo} surgical and behavioral procedures in head-fixed mice we direct the reader to existing resources (e.g. \cite{Guo2014FlowMice, Guo2014ProceduresMice} ). 

\subsubsection{Laser and power selection}
Zapit is compatible with any laser source with fast analog power modulation. Each setup will have unique experimental constraints but, as a guideline, we recommend using 1-2~mW time-averaged light power when using VGAT-ChR2 mice. We consistently see strong behavioral effects at these powers, and spatiotemporal resolution comparable with many other studies. Regular power measurements at the sample are critical to ensure stability and reproducibility. At 2~mW time-averaged power, and with expected light loss of around 50\% due to the skull and bone cement (\cite{Guo2014FlowMice}), heating of the neural tissue should be minimal (\cite{Stujenske2015ModelingOptogenetics}). Nevertheless, \textit{post-hoc} histological analysis is strongly recommended to confirm opsin expression and assess tissue integrity.

\subsubsection{Targeting precision}
With transgenic lines expressing opsin brain-wide (including VGAT-ChR2 mice), stimulation with high laser powers (e.g. 4~mW or above) may cause off-target behavioral effects due to light-scatter laterally and axially through cortex (\cite{Li2019SpatiotemporalCircuits,Babl2019TheInterneurons}). Ideally, laser parameters including power, duration and ramp-down should be informed by electrophysiological recordings and/or modeling of light propagation through biological medium (\cite{Yona2016RealisticApplications}). If limiting the impact of light dispersion is critical, opsins can be expressed via targeted viral injections. Under these conditions, expression can be restricted to highly localized circuits, genetically and anatomically defined subpopulations and sub-cellular compartments, and verified using \textit{post-hoc} histology.

\subsubsection{Sensory confounds}
Fast movements of galvo-mirrors generate auditory cues that can influence behavior, and we therefore recommend including moving the galvo-mirrors on all trials, including those where the laser power is 0~mW. The galvo drive waveforms are shaped to reduce noise emission, and the scanner unit can also be enclosed sound-attenuating housing to further mitigate these effects. To accurately interpret optogenetic behavioral effects, it is also critical to mask the mouse's ability to see the laser, whether through internal or external retinal stimulation \cite{Weiler}. We recommend masking laser-emitted light with color-matched LEDs placed above the mouse on every trial, matched to the stimulation profile of the laser (e.g., 40~Hz sinusoid). Extra caution should be taken when using longer wavelengths which easily propagate through brain tissue to the retina, and can be more difficult to mask (\cite{Danskin2015OptogeneticsArtifact}). Sensory confounds can also be addressed through repeated experiments and analyses in control animals that do not express the stimulated opsin (\cite{Gauld2025AActions}).

\subsubsection{Target region selection and interpretation} 
A major advantage of Zapit is the capacity to perform unbiased sampling of causal effects across the entire dorsal cortex. However, skull curvature and thickness should be considered, particularly in more lateral cortical regions. Differences in opsin expression due to tropisms in viral vectors or transgenic lines can also influence results and further experiments are required to determine whether observed effects are a direct result of regional inactivation, or derive from downstream brain regions (\cite{Wolff2018TheManipulations, Allen2017ApplicationSignaling}).

\subsection{Good practice: measures to report when using Zapit}
We encourage Zapit users to report a minimal set of standardized parameters to aid replication, comparability, and consistency across experiments and laboratories:

\begin{itemize}
    \item Average light power
    \item Peak light power
    \item Stimulation frequency and profile (e.g. 40~Hz sinusoid), and duty cycle
    \item Laser stimulation and ramp-down duration
    \item Color and frequency of masking light (if used)
\end{itemize}

\subsection{In Closing}
We hope Zapit is a tool that will grow and find uses outside of our institute and specialties. Potential new users are welcome to contact us for assistance in building or setting up the system. 
